\documentclass[11pt]{article}
\usepackage[utf8]{inputenc}
\usepackage{amsmath,amssymb,amsthm}
\usepackage{booktabs}
\usepackage[margin=2.5cm]{geometry}
\usepackage{hyperref}

\newtheorem{conjecture}{Conjecture}
\newtheorem{observation}{Observation}
\newtheorem{proposition}{Proposition}
\newtheorem{corollary}{Corollary}
\theoremstyle{definition}
\newtheorem{definition}{Definition}
\theoremstyle{remark}
\newtheorem{remark}{Remark}

\newcommand{\D}{\mathrm{D}}
\newcommand{\qbin}[2]{\left[\begin{smallmatrix}#1\\#2\end{smallmatrix}\right]}

\title{The defect of the differential expansion beyond the Alexander polynomial:\\
an experimental study on the Rolfsen--Hoste--Thistlethwaite table\\[2mm]
\large\emph{(draft)}}
\author{}
\date{}

\begin{document}
\maketitle

\begin{abstract}
Using colored HOMFLY polynomials $H_R(A,q)$ of all $2893$ knots with at most 12 crossings
and all representations $|R|\le 6$ (and symmetric representations up to $[10]$ for selected
arborescent knots) we study the factorization structure of the differential expansion.
We find that the defect is \emph{not} determined by the Alexander polynomial in general:
19 knots violate the rule $\delta=\deg\Delta-1$, and for some of them the defect is
fractional.  We propose to describe the defect by a tower of integers $e_m$, the $\lambda$-degrees
of the specialisations $P_m(q,\lambda)=H_{[r]}(A=q^{-m})|_{q^r=\lambda}$, whose $q\to1$ limit is
$\Delta(\lambda^2)^{m+1}$.  The tower is superadditive, bounded by $(m+1)\deg\Delta\le e_m\le
(m+1)g$, and the defect is its slope.  The former ``$\delta=-1$'' knots ($\Delta=1$) fit
naturally.  For non-rectangular representations we find a complete set of universal
reductions of $H_R$ on the lines $A=q^{-n}$, which replace the factorization of $H_R-1$.
\end{abstract}

\section{Conventions and data}

$H_R(A,q)$ is the reduced colored HOMFLY polynomial,
$A^{-1}H(L_+)-AH(L_-)=(q-q^{-1})H(L_0)$, $H_R(\text{unknot})=1$, topological framing.
We use
\[
  \D_n=Aq^n-A^{-1}q^{-n},\qquad [n]=\frac{q^n-q^{-n}}{q-q^{-1}},\qquad
  \qbin{r}{k}=\frac{[r]!}{[k]![r-k]!}.
\]
The data: all $R$ with $|R|\le 6$ for knots up to 11 crossings, $|R|\le 4$ (part of them $\le6$)
for 12 crossings; symmetric $[r]$, $r\le 10$, for the arborescent knots
$11n_{67},11n_{97},12n_{51},12n_{268},12n_{293},12n_{457}$ and $r\le 9$ for several
$\deg\Delta=2$ knots.  Methods: Hecke algebra (fundamental), tangle calculus with inclusive/exclusive
Racah matrices (arborescent knots; for $r\ge7$ the Racah matrices are evaluated numerically at the
interpolation points), cabling in the space of paths.  Every polynomial is reconstructed by dense
interpolation modulo one or two 31-bit primes and checked at fresh points modulo an independent prime
and at $q=1$; all anomalous cases below were cross-checked by an independent method (cabling of the
KnotInfo braid word), with exact agreement.

\section{Symmetric representations: the tower $e_m$}

\subsection{Extraction of $G_k$}
For symmetric representations we write
\begin{equation}\label{DE}
  H_{[r]}=1+\sum_{k=1}^{r}\qbin{r}{k}\;\D_{r+k-1}\cdots\D_{r}\;G_k(A,q)
\end{equation}
with $r$-independent $G_k$, extracted recursively from $H_{[k]}$.
\begin{observation}
For all 2893 knots and $k\le4$ (and $k\le 10$ where available) the division in (\ref{DE}) is exact,
i.e.\ all $G_k$ are Laurent polynomials (the ``plus'' structure is universal).
\end{observation}
The conventional defect-$\delta$ ansatz states
$G_k=F_k\,\D_{-1}\D_0\cdots\D_{m_k-2}$ with $m_k=\lfloor\frac{k-1}{\delta+1}\rfloor+1$ and
$\delta=\deg\Delta-1$.

\subsection{The tower}
\begin{observation}\label{obs:pers}
For all knots in the table:
(i) if $\D_n\,|\,G_k$ then $\D_n\,|\,G_{k'}$ for all $k'>k$ (persistence);
(ii) if $\D_n\,|\,G_k$ then $\D_{n-1}\,|\,G_k$ (consecutiveness);
(iii) every $\D_n$ divides $G_k$ at most once.
\end{observation}
Hence the whole structure is encoded in the integers
\[
  e_m+1:=\min\{k:\ \D_m\,|\,G_k\},\qquad m=-1,0,1,\dots
\]
Equivalently, on the line $A=q^{-m}$ the sum (\ref{DE}) terminates at $k=e_m$, and
\begin{definition}
$P_m(q,\lambda):=H_{[r]}(A=q^{-m},q)\big|_{q^r=\lambda}$ is a Laurent polynomial in $(q,\lambda)$
and $e_m=\tfrac12\deg_\lambda P_m$.
\end{definition}
\begin{itemize}
\item $m=-1$: $A=q$, $H_{[r]}=1$, so $e_{-1}=0$;
\item $m=0$: $A=1$, $P_0=\Delta(\lambda^2)$, so $e_0=\deg\Delta=:d$;
\item $m\ge1$: new invariants.  By $H_R(A,q)=H_{R^T}(A,-q^{-1})$ the line $A=q^{-m}$ corresponds to
$\mathfrak{sl}_m$ with the antisymmetric representations $[1^r]$, analytically continued in $r$
(equivalently, a supergroup $\mathfrak{gl}(1|m+1)$-type invariant with a continuous weight $\lambda$).
For $m=1$ we expect a relation to the Links--Gould invariant.
\end{itemize}

\begin{observation}[Melvin--Morton--Rozansky type limit]\label{obs:mmr}
$P_m(q=1,\lambda)=\Delta(\lambda^2)^{m+1}$
(checked for $4_1$ ($m=1,2$), $6_2$, $11n_{42}$, $11n_{67}$ ($m=1,2$), $12n_{293}$ ($m=1,2$)).
\end{observation}
Consequently $e_m\ge(m+1)d$.  For ``regular'' knots the leading $\lambda$-coefficient of $P_m$ is
$\pm1$, the leading coefficient of $\Delta^{m+1}$.  For anomalous knots $P_m$ has higher
$\lambda$-powers whose coefficients vanish at $q=1$, e.g.
\[
 11n_{67},\,m=1:\ (q^2-1)(q^2+1);\qquad
 11n_{42},\,m=1:\ -(q^2-1)(q^2+1)(q^4+q^2-1);\qquad
 12n_{293},\,m=2:\ \propto (q-q^{-1})^2[2][3].
\]
The Alexander polynomial is the $\hbar^0$ term of $P_m(e^\hbar,\lambda)$; anomalies come from
$O(\hbar)$ corrections of higher $\lambda$-degree.

\begin{observation}\label{obs:bounds}
For all knots, with $g$ the Seifert genus,
\[
  (m+1)\,d\ \le\ e_m\ \le\ (m+1)\,g ,
\]
and $f(n):=e_{n-1}$ is superadditive, $f(a)+f(b)\le f(a+b)$ (tested on 2217 knots with enough data).
\end{observation}
By Fekete's lemma the limit below exists, which we take as the definition of the defect:
\begin{equation}
  1+\delta:=\lim_{n\to\infty}\frac{f(n)}{n}=\sup_n\frac{f(n)}{n}\ \in[d,g].
\end{equation}
For $d=g$ (2857 of 2893 knots) $e_m=(m+1)d$ and $\delta=d-1$, the conventional defect.

\subsection{Anomalies}

\begin{table}[h]
\centering
\begin{tabular}{llcc l}
\toprule
knots & type & $d$ & $g$ & $e_{-1},e_0,e_1,e_2,\dots$\\
\midrule
$4_1$ (and all $d=g$) & regular & $1$ & $1$ & $0,1,2,3,4,5$\\
$6_2,\ 10_{129}$ & regular & $2$ & $2$ & $0,2,4,6,8$\\
$11n_{67},11n_{97},12n_{51},12n_{268}$ & I & $1$ & $2$ & $0,1,3,4,6,7,9$ \quad $(e_m=\lfloor\frac{3(m+1)}2\rfloor)$\\
$12n_{293},12n_{457}$ ($12n_{321}$, \dots) & II & $1$ & $2$ & $0,\mathbf 1,4,6,8$ \quad ($=2(m+1)$ for $m\ge1$)\\
$11n_{34},11n_{42},12n_{313},12n_{430}$ & $\Delta=1$ & $0$ & $3,2,2,2$ & $0,\mathbf 0,3,\ge5$\\
$12n_{750},12n_{830}$ & & $2$ & $3$ & $0,2,\ge6$\\
$11n_{73}$ & & $2$ & $3$ & $0,2,4,6,8$\\
$12n_{132}$ & & $2$ & $3$ & $0,2,4,\ge8$\\
\bottomrule
\end{tabular}
\caption{Towers $e_m$ (bold: the Alexander member falls below the slope of the tower).}
\end{table}

\begin{itemize}
\item All 19 knots violating $\delta=\deg\Delta-1$ have $g>\deg\Delta$, i.e.\ cancellations in knot
Floer homology; the converse fails ($11n_{73}$: $g=3$, regular slope 2).
\item Type II: only the Alexander member is degenerate ($e_0=1$ instead of 2); for $m\ge1$ the tower
saturates the genus bound, $e_m=(m+1)g$.
\item Type I: the slope is $3/2$, $\delta=\tfrac12$; on even $m$ one unit is lost (the sign
$(-1)^m$ of $A=(-1)^m q'^{\,m}$ after transposition).
\item $\Delta=1$ (former $\delta=-1$): $P_m(1,\lambda)=1$ for all $m$, all $\lambda$-dependence is in the
$\hbar$-corrections; $e_1=3$, so the old extrapolation of $e_0=0$ to the whole tower ($H\equiv1$) was
bound to fail.  The tower has slope $\ge 5/3$ and, since $11n_{34}$ and $11n_{42}$ are mutants
(identical symmetric colored HOMFLY, genera 3 and 2), the controlling invariant is
mutation-invariant: $1+\delta\le 2$.
\item $12n_{132}$: $e=0,2,4,\ge 8$ is not of the form $\lceil\cdot\rceil$ of a single slope at small
$m$; only the tower (and its asymptotic slope) is meaningful.
\end{itemize}

\begin{conjecture}
$e_m$ is the $\lambda$-degree of a $\mathfrak{gl}(1|m+1)$ (Links--Gould type) invariant $P_m(q,\lambda)$
with $P_m(1,\lambda)=\Delta(\lambda^2)^{m+1}$; $(m+1)\deg\Delta\le e_m\le (m+1)\,g$, and the defect
$1+\delta=\lim e_m/(m+1)$ can be fractional.  The integer rule $\delta=\deg\Delta-1$ holds exactly when the
$\hbar$-corrections to $P_m$ do not raise its $\lambda$-degree.
\end{conjecture}

\subsection{$P_m$ and the Links--Gould invariants}
Let $LG^{n,1}(t_0,t_1)$ be the Links--Gould invariant of $U_q\mathfrak{gl}(n|1)$ on the $2^n$-dimensional
typical module $V_\alpha$, with $t_0^{1/2}=q^{-\alpha}$, $t_1^{1/2}=q^{\alpha+1}$ for $n=2$ (Kohli--Tahar).
\begin{proposition}\label{prop:lg}
For every knot and $m\ge1$, $P_m(q,q^r)$ is, for each integer $r\ge m+1$, the $(1,1)$-tangle invariant of
$U_q\mathfrak{gl}(1|m+1)$ on $\mathrm{Sym}^r\mathbb{C}^{1|m+1}$, a $2^{m+1}$-dimensional typical module.  Hence
$P_m$ is $LG^{m+1,1}$ with the continuous weight $\alpha=r$, i.e.\ $q^r=\lambda$; for $m=1$
\begin{equation}\label{LG}
  P_1(q,\lambda)\;=\;LG^{2,1}\big(t_0=\lambda^{-2},\,t_1=q^2\lambda^{2}\big)\qquad(\text{up to mirror}).
\end{equation}
\end{proposition}
\begin{proof}[Proof (sketch)]
(1) \emph{Hecke algebra.}  Write the $R$-cabled $(1,1)$-tangle of $K$ as an element of $H_{n|R|}(q)$ and apply
Ocneanu's conditional expectation to the last $(n-1)|R|$ strands and the idempotent $e_R$; since $e_RH_{|R|}e_R=
\mathbb{C}e_R$ the result is $H_R(A,q)\,e_R$.  The expectation is a polynomial in $A^{\pm1}$ and
$\frac{A-A^{-1}}{q-q^{-1}}$, so this holds at $A=q^{-m}$, where $\dim_qR$ vanishes for $r\ge2$ and the closed
(unreduced) invariant carries no information.
(2) \emph{Super Schur--Weyl.}  For $V=\mathbb{C}^{1|m+1}$ the graded flip deformed to the Perk--Schultz matrix gives
a representation of $H_N(q)$ on $V^{\otimes N}$, the $U_q\mathfrak{gl}(1|m+1)$ braiding, under which the Markov
expectation with $A=q^{-m}$ becomes the partial quantum supertrace (both are the unique Markov expectations
with these parameters) and $e_{[r]}$ projects onto $\mathrm{Sym}^r V$.  So $H_{[r]}(q^{-m},q)$ is the scalar by
which the $(1,1)$-tangle acts on $\mathrm{Sym}^rV$.
(3) \emph{The module.}  $\mathrm{Sym}^rV$ has basis $e_0^{r-|S|}\prod_{i\in S}e_i$, $S\subseteq\{1,\dots,m+1\}$,
so it has dimension $2^{m+1}$ for $r\ge m+1$, highest weight $(r\,|\,0,\dots,0)$, and it is typical.  Under
$\mathfrak{gl}(1|m+1)\cong\mathfrak{gl}(m+1|1)$ (parity change, which mirrors the invariant) it is the module
$V_\alpha$ with $\alpha=r$, up to a one-dimensional twist which only changes the framing factor.
(4) Both sides are Laurent polynomials in $q^\alpha$ (resp.\ $\lambda$) and agree at $\lambda=q^r$ for all
$r\ge m+1$, hence identically.
\end{proof}
Steps (1)--(2) were checked independently of the Hecke/Racah methods that produced the data: a vertex model
built from the Perk--Schultz matrix, fused with the $q$-symmetrizer and closed with the pivotal element
(\texttt{scripts/supergroup\_vertex\_check.py}) reproduces $H_{[r]}(q^{-m},q)$ (in the mirror convention) at
random real $q$: for $\mathfrak{gl}(1|2)$, $r=2,3$, on $255$ knots (all knots up to $10$ crossings of braid index
$\le5$, and those of $11n_{34},11n_{42},11n_{67},11n_{73},11n_{97},12n_{23},12n_{51},12n_{63},12n_{268},12n_{293},
12n_{313},12n_{430},12n_{457}$ with braid index $\le5$), for $r=4$ on $3_1,4_1,5_1,5_2,7_1$, and for
$\mathfrak{gl}(1|3)$, $r=3$ (dimension $8$, i.e.\ $P_2$), on the $62$ knots above of braid index $\le3$.
The value of $LG^{2,1}$ itself was not computed independently; the dictionary $\alpha=r$ agrees with the
specialisations below and, for the trefoil, (\ref{LG}) gives $1-t_0-t_1+t_0^2+2t_0t_1+t_1^2-t_0^2t_1-t_0t_1^2$.
The properties of $LG^{2,1}$ were also checked on $P_1$ over the table
(\texttt{scripts/links\_gould\_check.py}, \texttt{data/links\_gould\_p1\_scan.tsv}).
\begin{observation}\label{obs:lg}
For all $1068$ knots whose data determine $P_1$ or its top part (302 of them completely):
(i) $P_1(q,\lambda)=P_1(q,q/\lambda)$ ($t_0\leftrightarrow t_1$ symmetry);
(ii) $P_1(1,\lambda)=\Delta(\lambda^2)^2$ ($LG(t_0,t_0^{-1})=\Delta(t_0)^2$);
(iii) $P_1(i,\lambda)=\pm\Delta(\lambda^4)$ ($LG(t_0,-t_0^{-1})=\Delta(t_0^2)$);
(iv) the mutants $11n_{34}$, $11n_{42}$ (Conway, Kinoshita--Terasaka) have equal $P_1$, as they have equal $LG$.
\end{observation}
Under (\ref{LG}) the $(t_0,t_1)$-span used by Kohli and Tahar (max minus min of $a-b$ over the monomials
$t_0^at_1^b$) equals $\tfrac12\,\mathrm{span}_\lambda P_1=2e_1$.  Their theorem
$\mathrm{span}\,LG\le 4g$, with equality for alternating knots, therefore \emph{is} the bound $e_1\le 2g$ of
Observation~\ref{obs:bounds}, with $e_1=2g$ for alternating knots; both hold on the whole table.
Of the 302 knots with complete $P_1$, $288$ have $e_1=2g$; the $14$ with $e_1<2g$ are
$11n_{34},11n_{42},11n_{67},11n_{73},11n_{97},12n_{51},12n_{56},12n_{63},12n_{87},12n_{132},12n_{221},
12n_{268},12n_{313},12n_{430}$.

\begin{observation}[roots of unity]\label{obs:roots}
Let $q$ be a primitive $2k$-th root of unity, $2\le k\le m+1$, and $m+1=a+kb$, $b=\lfloor\frac{m+1}k\rfloor$.  Then
\[
  P_m(q,\lambda)=\pm\,\Delta(\lambda^2)^{a}\,\Delta(\lambda^{2k})^{b}
\]
(checked for $m=1$ on all 302 knots with complete $P_1$, for $m=2$ on the 27 knots with complete $P_2$, and for $m=3$ on $3_1$; it fails for $k=m+2$).
\end{observation}
This should follow from $H_{[r+k]}=H_{[r]}H_{[k]}$ at $q^{2k}=1$ together with the $q$--$A$ duality of the
Kashaev specialisation (arXiv:1505.06170).  It forces every coefficient of $\lambda^j$, $j>2(m+1)d$, in
$P_m$ to be divisible by $\prod_{k=1}^{m+1}\Phi_{2k}(q)$ and by $\Phi_1(q)$, which is the origin of the
factors in the anomalous coefficients above: $(q^2-1)(q^2+1)=\Phi_1\Phi_2\Phi_4$ for $m=1$, and
$(q-q^{-1})^2[2][3]\ni\Phi_1\Phi_2\Phi_4\Phi_6$ for $m=2$.

\paragraph{The slope is not decided by $\tau$ or $\sigma$.}  For $g=2$, $d=1$ the Type I knots
($e_1=3$) all have $\tau=\sigma=0$ and $12n_{293},12n_{321},12n_{411},12n_{457},12n_{124},12n_{519}$
($e_1=4$) have $(\tau,\sigma)\neq(0,0)$, but $12n_{23}$ has $\tau=\sigma=0$ and $e_1=4$, and for $g=3$
$12n_{63}$ ($\tau=-1,\sigma=2$) has $e_1=4<2g$.  In the examples we compared, the top two Alexander gradings of $\widehat{HFK}$ do not
separate the two types either.

\section{Non-rectangular representations}

For rectangular $R=[r^s]$, $H_R-1$ is divisible exactly by $\D_r\D_{-s}$
($R$ is trivial for $\mathfrak{sl}_s$, and for $\mathfrak{sl}_r$ after transposition).
For non-rectangular $R$, $H_R-1$ is divisible by no $\D_n$ for every knot.
The role of these factors is taken over by universal reductions.

\begin{observation}[universal reductions]\label{obs:red}
On the line $A=q^{-n}$, $n\ge0$ (for $n<0$ apply the rule to $R^T$ and $-n$):
\begin{enumerate}
\item[(a)] if $n\ge R_1$: $H_R=H_{R^*}$, $R^*=\big(\text{complement of }R^T\text{ in the }n\times\ell(R)
\text{ box}\big)^T$ (conjugation in $\mathfrak{sl}_n$);
\item[(b)] if $0\le n<R_1$ and $R_2=n+1$: $H_R=H_{(R_1+1,R_3,R_4,\dots)}$;
\item[(c)] otherwise there is no universal relation.
\end{enumerate}
Two representations coincide on the line iff their normal forms (minimal elements of the reduction
orbits) coincide.  Rule (b) for $n=0$ is the hook property $H_R(A=1)=\Delta(q^{2|R|})$; for $n\ge1$ it
gives e.g.\ $[2,2]\equiv[3]$, $[3,2]\equiv[4]$, $[3,2,1]\equiv[4,1]$ at $A=q^{-1}$,
$[3,3]\equiv[4]$ at $A=q^{-2}$, and chains $[2,2,2]\to[3,2]\to[4]$.
\end{observation}
Verification: all knots, $|R|\le6$, $|n|\le8$, $9.5\cdot10^6$ pairs $(R,R')$: no predicted identity fails,
and no further pairwise identity holds except at $A=1$ for the four knots with $\Delta=1$.
No multi-term linear relation $H_R=\sum_{R'}c_{R'}(q)H_{R'}$ with knot-independent coefficients exists
beyond these (apart from consequences of the Alexander property at $A=1$).

\paragraph{Knot-dependent structure.}
On the lines where (c) applies, linear relations with knot-independent coefficients appear only within
classes of knots.  For $d=g=1$ (defect 0) at $A=q^{-1}$:
\[
  H_{[4]}\in\mathrm{span}\{1,H_{[2]},H_{[3]}\},\qquad
  H_{[4,1]}\in\mathrm{span}\{1,H_{[2,1]},H_{[3,1]}\},
\]
while $H_{[3,1]}\notin\mathrm{span}\{1,H_{[2,1]}\}$: the hook family $[r,1]$ has the same tower member
$e_1=2$ as the symmetric family, with different coefficient functions.  This suggests the general
formulation: the differential expansion of an arbitrary $R$ is organised by families $R(\lambda)$
(e.g.\ $[r]$, $[r,1]$, $[r,2]$, $[r,r-1]$) on the lines not covered by (a), (b), with towers $e_m$ per
family.

\section{Open questions}
\begin{enumerate}
\item Fill in the proof of Proposition~\ref{prop:lg} (the Markov-expectation argument in (2), and the
one-dimensional twist in (3)), and compare (\ref{LG}) with an independent evaluation of $LG^{2,1}$.
Extend the Kohli--Tahar bound to $LG^{n,1}$: by Proposition~\ref{prop:lg} a bound
$\mathrm{span}\,LG^{n,1}\le 2ng$ (span in $q^{2\alpha}$, the normalisation in which their theorem reads
$\mathrm{span}\,LG^{2,1}\le4g$) is exactly $e_{n-1}\le ng$.
\item What selects the slope inside $[d,g]$?  Types I and II have equal $d,g$ and similar knot Floer
homology, but slopes $3/2$ and $2$.
\item Towers of the hook and two-row families for $|R|\le 10$ (arborescent knots).
\item A conceptual derivation of rule (b) (supergroup / Kazhdan--Lusztig type).
\end{enumerate}

\end{document}
